\documentclass[../../../include/open-logic-section]{subfiles}

\begin{document}

\olfileid[pt]{sth}{spine}{foundation}

\olsection{Fundação}

Estamos apenas \emph{quase} concluídos---e não \emph{totalmente} concluídos---porque nada em
$\ZFminus$ garante que \emph{todo} o conjunto está em algum $V_\alpha$,
O.s.\@, que todo o conjunto é formado em algum estágio. 

Há, contudo, um sentido (matemático) bastante direto em que não
\emph{importa} se existem conjuntos fora da hierarquia. (Se os houver,
podemos simplesmente ignorá-los.) Mas motivamos o nosso \emph{conceito} de conjunto com a ideia de que todo o conjunto é formado em
algum estágio (ver \stageshier{} na \olref[z][story]{sec}). Logo
queremos excluir a possibilidade de conjuntos que caem fora da
hierarquia. Por conseguinte, devemos acrescentar um novo axioma, que assegura
que todo o conjunto surge em algum lugar da hierarquia. 

Uma vez que os $V_\alpha$ são os nossos estágios, poderíamos simplesmente considerar acrescentar o seguinte como axioma:

\begin{defish}
\emph{Regularidade.} $\forall A \exists \alpha\, A \subseteq V_\alpha$
\end{defish}

Seria uma abordagem perfeitamente razoável. Contudo, por razões
que se explicarão adiante, adotaremos antes um axioma alternativo:

\begin{axiom}[Fundação]
$(\forall A \neq \emptyset)(\exists B \in A)A \cap B = \emptyset$.
\end{axiom}

Com algum trabalho, podemos mostrar (em $\ZFminus$) que a Fundação implica a Regularidade:
\begin{defn}
Para cada conjunto $A$, definimos:
\begin{align*}
	\text{cl}_0(A) &= A,\\
	\text{cl}_{n+1}(A) &= \bigcup \text{cl}_n(A),\\
	\trcl{A} &= \bigcup_{n < \omega} \text{cl}_{n}(A).
\end{align*}
Chamamos a $\trcl{A}$ o \emph{fecho transitivo} de $A$.
\end{defn}
\noindent
O nome «fecho transitivo» é adequado:
\begin{prop}\ollabel{subsetoftrcl}
$A \subseteq \trcl{A}$ e $\trcl{A}$ é um conjunto transitivo. 
\end{prop}

\begin{proof}
Evidentemente $A = \text{cl}_0(A) \subseteq \trcl{A}$. E se $x
\in b \in \trcl{A}$, então $b \in \text{cl}_n(A)$ para algum $n$, logo $x
\in \text{cl}_{n+1}(A) \subseteq \trcl{A}$. 
\end{proof}

\begin{lem}\ollabel{lem:TransitiveWellFounded}
Se $A$ é um conjunto transitivo, então existe algum $\alpha$ tal que $A
\subseteq V_\alpha$.
\end{lem}

\begin{proof}
Recordando a definição de «$\supstrict(X)$» em
\olref[ordinals][opps]{defsupstrict}, definimos dois conjuntos:
\begin{align*}
	D  &= \Setabs{x \in A}{\forall \delta\ x \nsubseteq V_\delta}\\
	\alpha &= \supstrict\Setabs{\delta}{(\exists x \in A)
	(x \subseteq V_\delta \land (\forall \gamma \in \delta)x \nsubseteq V_\gamma)}
\end{align*}
Suponhamos $D = \emptyset$. Se $x \in A$, então existe algum $\delta$ tal que $x \subseteq V_\delta$ e, pela boa ordenação dos ordinais, $(\forall \gamma \in \delta)x \nsubseteq V_\gamma$; logo $\delta \in \alpha$ e assim $x \in V_\alpha$ por
\olref[spine][Valphabasic]{Valphabasicprops}. Logo $A \subseteq
V_{\alpha}$, como pretendido. 

Basta pois mostrar que $D = \emptyset$. Por redução ao absurdo, suponhamos
o contrário. Pela Fundação, existe algum $B \in D \subseteq A$ tal que $D \cap B
= \emptyset$. Se $x \in B$ então $x \in A$, pois $A$ é transitivo,
e como $x \notin D$, segue que $\exists \delta\ x \subseteq
V_\delta$. Seja agora
\[
	\beta = \supstrict\Setabs{\delta}{(\exists x \in B)(x \subseteq V_\delta \land (\forall \gamma < \delta)x \nsubseteq V_\gamma)}.
\]
Como antes, $B \subseteq V_\beta$, contradizendo a afirmação de que $B \in
D$.	
\end{proof}

\begin{thm}\ollabel{zfentailsregularity}
A Regularidade vale.
\end{thm}

\begin{proof}
Fixemos $A$; agora $A \subseteq \trcl{A}$ por \olref{subsetoftrcl}, que é transitivo. Logo existe algum $\alpha$ tal que $A \subseteq \trcl{A} \subseteq V_\alpha$ por \olref{lem:TransitiveWellFounded}
\end{proof}

Estes resultados mostram que $\ZFminus$ prova a implicação
$\emph{Fundação}\Rightarrow\emph{Regularidade}$. Em
\olref[spine][rank]{zfminusregularityfoundation}, mostraremos que
$\ZFminus$ prova $\emph{Regularidade}\Rightarrow\emph{Fundação}$. Como
tal, a Fundação e a Regularidade são \emph{equivalentes} (módulo
$\ZFminus$). Mas isto significa que, dado $\ZFminus$, podemos justificar a
Fundação notando que é equivalente à Regularidade. E podemos
justificar a Regularidade de imediato com base em \stageshier{}. 

\end{document}
